\section{Introduction}

Local distortion of periodic-point counts often changes a rational
dynamical zeta function into one with a natural boundary. When several
archimedean phases have the same maximal modulus, however, a
natural-boundary argument must account for cancellation between them.
This issue persists even if every individual phase-weighted summand
has a natural boundary: the sum need not have one.
We resolve this cancellation question for finite-adelic radial
coefficients with periodic data by an exact finite test.

The coefficients under consideration have the form
\begin{equation}\label{eq:intro-coeff}
 a_n=\left(\sum_{j=1}^{J}c_j\eta_j^n\right)r_n
 \prod_{p\in S}p^{-s_{p,n}v_p(n)-t_{p,n}p^{v_p(n)}}.
\end{equation}
Here \(S\) and the phase set are finite, \(|\eta_j|=1\), the sequence
\(r_n\) is periodic and complex, and all \(s_{p,n},t_{p,n}\) are
periodic nonnegative real numbers. No algebraicity, height bound,
positivity of the phase weights, or independence condition on the
phases is imposed. Theorem~\ref{thm:main} states that
\(\sum_{n\geq1}a_nz^n\) is rational precisely when every potentially
unbounded active valuation fibre is annihilated by the appropriate
grouped phase weight. Otherwise every point of the unit circle
obstructs meromorphic continuation.

The test is finite because only the common residue period, the finite
phase classes modulo roots of unity, and the primes in \(S\) enter it.
Its proof is not a no-cancellation principle for arbitrary functions.
We first construct an absolutely summable Fourier expansion on the
closure of the diagonal integers in a finite residue group times
finitely many \(p\)-adic integer groups. We then combine all equal
frequencies. An active fibre forces unbounded conductor in this
actual expansion. Local radial invariance produces full root grids
along an unbounded sequence of conductors in the nonzero Fourier support,
with the same nonzero coefficient throughout each grid.
Thus the \emph{actual} atomic measure, rather
than a formal decomposition of it, has dense nonzero atoms.

Natural-boundary questions for algebraic dynamics have a substantial
history. Bell--Miles--Ward developed a P\'olya--Carlson programme for
algebraic systems, including finite-place distortion and the role of
archimedean unit roots \citep{bmw2014}. Byszewski--Cornelissen treated
abelian varieties in positive characteristic, with a valuation-kernel
dichotomy in the appendix by Royals--Ward \citep{bc2018}.
Byszewski--Cornelissen--Houben organized a much wider class through
finite-adelically distorted systems \citep{bch2024}.
Their inspected Theorem~11.3.8 assumes a unique dominant root and
integral wild data. Our dynamical application removes these two
assumptions for the finite-adelic presentation considered here.

It is important to distinguish this conclusion from known no-wild
results. Baril Boudreau--Holmes--Nguyen prove a natural-boundary
alternative for stable rational series perturbed by almost
quasi-polynomials \citep{bhn2025}. Their theorem already gives the
no-wild nonhyperbolic finite-adelic conclusion; we give the deduction
in Section~\ref{sec:dynamics}. The genuinely wild example in
Section~\ref{sec:example} fails their sublinear-height hypothesis.
Likewise, the normalized example does not satisfy the all-embeddings
unit-disc convergence condition in the criterion of
Bell--Gunn--Nguyen--Saunders \citep{bgns2023}.
These failed sufficient hypotheses delimit the applications of those
theorems; they are not, by themselves, priority assertions.

\begin{table}[t]
\centering
\small
\begin{tabular}{@{}>{\raggedright\arraybackslash}p{0.23\textwidth}
>{\raggedright\arraybackslash}p{0.35\textwidth}
>{\raggedright\arraybackslash}p{0.34\textwidth}@{}}
\toprule
Result & Relevant hypotheses & Role here\\
\midrule
\citep[Theorem 11.3.8]{bch2024}
& Finite-adelic dynamics; unique dominant root; integral wild data
& Classical distorted natural-boundary theorem\\[3pt]
\citep[Theorem 2.14]{bhn2025}
& Stable rational series; almost quasi-polynomial perturbation of
sublinear height
& Already implies the no-wild application\\[3pt]
Theorem~\ref{thm:main}
& Finite phases and primes; periodic nonnegative real exponents;
complex weights
& Exact active-fibre test, including phase masking and wild data\\
\bottomrule
\end{tabular}
\caption{Different hypothesis classes in the comparison. The theorem
locators in the first two rows refer to the public versions specified
in the text and bibliography.}
\label{tab:scope}
\end{table}

The Fourier-to-boundary passage is classical in character.
For example, Knill--Lesieutre use Fourier expansions and radial
singularities for almost-periodic coefficients
\citep[Proposition 2.2 in the revised manuscript]{kl2012}.
Cornelissen--Park also use Fourier analysis of finite-adelic
coefficients, including archimedean phases, to study orbit-decomposition
statistics \citep{cp2026}. We do not claim a new general Cauchy-transform
principle or the first Fourier treatment of these coefficients.
The point is the exact active-fibre criterion and the conductor-grid
argument proving dense \emph{noncancelled} masses in its active branch.

For clarity about source versions, our detailed comparisons use
\citep{bch2024} in its 19 April 2024 public version and
\citep{bhn2025} in its 16 July 2023 preprint version.
The latter's journal metadata has been verified, but its final text
has not been compared line by line with that preprint; no assertion
that a later book or final revision is unchanged is needed for the
proofs below.
